\section{Method}
\label{sec:method}

Cross-frequency coupling is a relation between two frequency bands, and a
token grid built from per-band content has no cell in which to write a
relation. CroFreMo adds one. It maps a multichannel EEG window to a
three-axis token grid of electrodes $\times$ frequency bands $\times$ time
patches in which every band contributes two token rows: one carrying the
band waveform, and one carrying an amplitude-preserving \emph{interaction}
representation derived from the band's measured phase coupling to slower
bands. A factorized transformer mixes the grid along each axis in turn.
Two principles govern the design: coupling may add \emph{structure} but never
\emph{energy}; and every coupling pathway is gated, so that the
\emph{uncoupled} model remains an exactly recoverable configuration and
coupling contributes only when supported by optimization. The full model has
1.6M parameters.

\subsection{Cross-frequency coupling as a signal phenomenon}
\label{sec:phenomenon}

Let $x \in \mathbb{R}^{C \times T}$ be a window of $C$ electrodes at
sampling rate $f_s$. For a band-limited component $x_b(t)$, the analytic
signal $z_b = x_b + i\,\mathcal{H}[x_b]$ splits the band into an amplitude
envelope $A_b(t) = |z_b(t)|$ and a unit phase $p_b(t) = z_b(t)/|z_b(t)|
= e^{i\phi_b(t)}$. Phase--amplitude coupling (PAC) between a slow band $i$
and a fast band $j$ is the statement that $A_j(t)$ is not independent of
$\phi_i(t)$: the fast envelope is systematically larger at some phases of
the slow rhythm than at others. Its classical estimators are all functionals
of the joint distribution of $(\phi_i, A_j)$. The modulation index of
\citet{tort2010measuring} bins $\phi_i$, averages $A_j$ within each bin, and
measures the Kullback--Leibler divergence of the resulting distribution
from uniform; the mean-vector length of \citet{canolty2006high} is the
first circular moment,
\begin{equation}
\mathrm{MVL}_{ij} \;=\; \Bigl|\,\frac{1}{T}\sum_{t} A_j(t)\,e^{i\phi_i(t)}\Bigr| ,
\label{eq:mvl-classical}
\end{equation}
whose magnitude is zero when $A_j$ carries no phase information and grows
with the depth of modulation, and whose argument is the preferred phase.
Both are computed over whole recordings, with band pairs fixed in advance,
as a single scalar per pair.

CroFreMo keeps the first circular moment but changes three things about how
it is computed. A learnable sinc filterbank \citep{ravanelli2018sincnet}
with $n_b$ band-pass filters, parameterized by cutoff pairs initialized on
a logarithmic (constant-$Q$) grid, produces the band signals $x_b(t)$, so
the bands are learned rather than fixed. Time is cut into $P$
non-overlapping patches of length $\ell$ (1\,s throughout), and the moment is
taken within each patch and channel rather than over the recording, so
that transient coupling events survive at patch resolution. And the
amplitude is centered before the moment is taken,
\begin{equation}
Z_{ij} \;=\; \frac{1}{\ell}\sum_{t \in \text{patch}}
p_i(t)\,\bigl(A_j(t) - \bar A_j\bigr),
\label{eq:mvl}
\end{equation}
which removes the term $\bar A_j \cdot \frac{1}{\ell}\sum_t p_i(t)$ that a
non-uniform phase distribution alone would contribute, so that $Z_{ij}$
responds to the \emph{co-variation} of envelope and phase and not to
the mean envelope. $Z_{ij}$ is kept complex: $|Z_{ij}|$ is the strength of
band-$i$ phase coupling onto band-$j$ amplitude, $\angle Z_{ij}$ its
preferred phase. For $n_b$ bands this yields a lower-triangular
$n_b \times n_b$ complex matrix per (channel, patch), computed by the
tokenizer for every window it sees.

\subsection{Coupling as token content: the duplex tokenizer}
\label{sec:tokenizer}

Each (channel, band, patch) yields three linear projections: a waveform
token $\mathbf{r}_b = W_r\,x_b[\text{patch}] \in \mathbb{R}^d$, an
amplitude feature
$\mathbf{a}_b = L_a(A_b[\text{patch}]) \in \mathbb{R}^{d/2}$, and a complex
phase feature
$\tilde{\mathbf{p}}_b = L_p(p_b[\text{patch}]) \in \mathbb{C}^{d/2}$, where
$L_p$ is linear and bias-free. The waveform token is what a per-band
tokenizer would produce on its own; the two analytic features are the raw
material for the interaction token, and the three properties below say how
that token must be built so that it carries coupling and nothing else.

\begin{proposition}[Phase-reference equivariance]
\label{prop:equivariance}
Let band $i$'s phase convention be shifted by a constant $\delta_i$, i.e.\
$p_i \mapsto e^{i\delta_i} p_i$. Then $\tilde{\mathbf{p}}_i \mapsto
e^{i\delta_i}\tilde{\mathbf{p}}_i$ and $Z_{ij} \mapsto e^{i\delta_i} Z_{ij}$
for every $j$.
\end{proposition}
\begin{proof}
$L_p$ is $\mathbb{C}$-linear without bias, so $L_p(e^{i\delta}p) =
e^{i\delta}L_p(p)$; and $e^{i\delta_i}$ factors out of the sum in
Eq.~(\ref{eq:mvl}).
\end{proof}

Raw phase depends on an arbitrary reference (the Hilbert transform fixes
it only up to the analysis convention), so by
Proposition~\ref{prop:equivariance} the phase feature
$\tilde{\mathbf{p}}_b$ is not, by itself, a well-defined token content.
We therefore let coupling enter through the \emph{aligned phase}, which
combines slower bands' phase features under the measured coupling geometry:
\begin{equation}
\mathbf{u}_j \;=\; \sum_{i<j} \alpha_{ij}\, e^{-i \angle Z_{ij}}\,
\tilde{\mathbf{p}}_i,
\qquad
\alpha_{ij} = \frac{|Z_{ij}|}{\sum_{i'<j}|Z_{i'j}|},
\label{eq:aligned}
\end{equation}
with $\mathbf{u}_0 = \tilde{\mathbf{p}}_0$ for the slowest band.

\begin{proposition}[Phase-reference invariance]
\label{prop:invariance}
$\mathbf{u}_j$ is invariant to the phase convention of every band.
\end{proposition}
\begin{proof}
Under $p_i \mapsto e^{i\delta_i} p_i$, Proposition~\ref{prop:equivariance}
gives $\tilde{\mathbf{p}}_i \mapsto e^{i\delta_i}\tilde{\mathbf{p}}_i$ and
$e^{-i\angle Z_{ij}} \mapsto e^{-i\delta_i}e^{-i\angle Z_{ij}}$; the two
factors cancel term by term in Eq.~(\ref{eq:aligned}), and $\alpha_{ij}$
depends on $|Z_{ij}|$ only.
\end{proof}

Proposition~\ref{prop:invariance} is the reason coupling enters through
Eq.~(\ref{eq:aligned}) rather than through the raw phases: the rotation by
$e^{-i\angle Z_{ij}}$ re-expresses each slow band's phase feature relative
to the phase at which it actually modulates band $j$, which is a physical
quantity, and the weights $\alpha_{ij}$ let the strongest modulators of
band $j$ dominate its aligned phase.

\begin{proposition}[Magnitude preservation]
\label{prop:isometry}
Define the interaction token
\begin{equation}
\mathbf{h}_j \;=\; \mathbf{a}_j \odot
\frac{\mathbf{u}_j}{|\mathbf{u}_j|} .
\label{eq:rotation}
\end{equation}
Then $|\mathbf{h}_j| = |\mathbf{a}_j|$ elementwise, for every value of the
coupling.
\end{proposition}
\begin{proof}
Each component of $\mathbf{u}_j/|\mathbf{u}_j|$ has unit modulus.
\end{proof}

Multiplying by $\mathbf{u}_j$ directly would rescale the amplitude feature
by $|\mathbf{u}_j|$, confounding coupling geometry with changes in band
power, which is the one quantity every EEG model already exploits. Under
Eq.~(\ref{eq:rotation}) band power reaches the encoder untouched; the
coupling geometry is encoded purely in phase, and coupling \emph{strength}
still enters through the weights $\alpha_{ij}$ of Eq.~(\ref{eq:aligned}).
The isometry also makes the ablation interpretable: any gain over the
uncoupled baseline cannot be an energy artifact.

\paragraph{The duplex grid.}
The final grid stacks, for each physical band $b$, a \emph{fused} row and
a gated \emph{interaction} row:
\begin{equation}
\text{row}_b = \mathbf{r}_b + \boldsymbol{\beta}_b \odot \mathbf{h}_b, \qquad
\text{row}_{n_b + b} = \gamma_b\, \mathbf{h}_b,
\label{eq:duplex}
\end{equation}
with $\boldsymbol{\beta}_b \in \mathbb{R}^d$ initialized to zero and scalar
$\gamma_b$ initialized to one. The two rows give coupling two different
routes into the encoder. The fused row lets coupling modulate the band's
own waveform representation in place; because $\boldsymbol{\beta}_b$
starts at zero, its gradient at initialization is
$\partial\mathcal{L}/\partial\boldsymbol{\beta}_b =
\partial\mathcal{L}/\partial\,\text{row}_b \odot \mathbf{h}_b$, so
coupling enters this route only along the directions in which the
interaction token is correlated with the loss gradient, and the fused rows
coincide with the uncoupled model until it does. The interaction row makes
coupling a separately attendable feature from the first step, which is what
event-level tasks turn out to need (Section~\ref{sec:results}).

\begin{proposition}[Exact recoverability]
\label{prop:recover}
With a mean-pooled readout followed by a normalization layer, the
configuration $\boldsymbol{\beta}_b = 0$, $\gamma_b = 0$ for all $b$
computes the same function as the uncoupled model with $n_b$ waveform rows.
\end{proposition}
\begin{proof}
At $\boldsymbol{\beta} = 0$ the fused rows are the waveform tokens. At
$\gamma = 0$ the interaction rows are zero vectors; self-attention with
zero rows present differs from self-attention without them only through
those rows' own outputs, and mean-pooling over $2n_b$ rows of which $n_b$
are zero rescales the pooled representation by $1/2$, which the head's
normalization removes.
\end{proof}

The uncoupled model is thus a reachable point of the parameterization
rather than a separate architecture, and the trained values of
$\boldsymbol{\beta}$ and $\gamma$ are reported per corpus as measurements
of how much coupling each task retained.

\subsection{Tri-axial encoder}
\label{sec:encoder}

The grid $(C \times 2n_b \times P \times d)$ is processed by $N$ pre-norm
blocks, each applying self-attention along one axis at a time --- time
(with rotary position encoding), electrodes (with a learned per-electrode
embedding), and frequency rows (with an embedding of each band's learned
center frequency and bandwidth) --- followed by a feed-forward layer, each
as a residual. Factorizing by axis keeps attention quadratic in one short
axis at a time and gives each axis its natural positional structure. All
coupling information reaches the encoder through the token content itself;
the frequency sub-layer is plain attention over the $2n_b$ rows. A
mean-pooled linear head produces predictions. We use $n_b{=}8$, $d{=}128$,
$N{=}6$, four heads.

\subsection{A pretraining objective for coupling}
\label{sec:objective}

If cross-frequency structure belongs in the token, it is natural to ask
whether it also belongs in the pretraining target. We construct the
objective that puts it there; Section~\ref{sec:results} evaluates it
empirically, to convergence.

\paragraph{Masking.} A random subset $\mathcal{M}$ of \emph{physical}
bands is masked per channel. Because the interaction row of band $b$
carries $\mathbf{a}_b$ by construction (Eq.~(\ref{eq:rotation})), masking
hides \emph{both} of $b$'s rows --- fused and interaction --- or the
visible interaction row would hand the reconstruction its own target; both
are replaced by a shared learned mask token before the positional
embeddings. Entries of the coupling matrix consumed by the encoder are
zeroed wherever either band is masked, so the target cannot travel through
the coupling input.

\paragraph{Objective.} The loss combines an amplitude term and a coupling
term,
\begin{equation}
\mathcal{L} \;=\; \mathcal{L}_{\mathrm{amp}}
\;+\; \lambda\, \mathcal{L}_{\mathrm{cfc}},
\qquad \lambda = 1.
\label{eq:objective}
\end{equation}
$\mathcal{L}_{\mathrm{amp}}$ asks the model to predict, from each masked
band's fused-row output, the patch log-amplitude standardized \emph{within
each band} across the batch --- the rebalancing of
\citet{yu2026understanding}, adopted as-is --- so that every band's
structure is rewarded equally rather than in proportion to its
$1/f$-dominated variance:
\begin{equation}
\mathcal{L}_{\mathrm{amp}} = \frac{1}{|\mathcal{M}|}
\sum_{j \in \mathcal{M}}
\bigl\| \hat{y}_j - \mathrm{std}_b\!\bigl(\log \bar A_j\bigr) \bigr\|^2 .
\label{eq:lamp}
\end{equation}
$\mathcal{L}_{\mathrm{cfc}}$ asks it to predict the masked band's coupling
column --- the real and imaginary parts of $Z_{ij}$ --- restricted to
\emph{visible} phase bands $i$, since the phase required to define
$Z_{ij}$ is unavailable when band $i$ is masked:
\begin{equation}
\mathcal{L}_{\mathrm{cfc}} = \frac{1}{Z_{\mathrm{sup}}}
\sum_{j \in \mathcal{M}} \sum_{i \notin \mathcal{M}}
\bigl\| \hat{Z}_{ij} - Z_{ij} \bigr\|^2,
\label{eq:lcfc}
\end{equation}
with $Z_{\mathrm{sup}}$ the count of supervised entries. Pretraining uses
the same corpora pool for every objective variant compared in
Section~\ref{sec:results}; per-variant schedules are given in the appendix.
